% %File: formatting-instructions-latex-2026.tex
% %release 2026.0
% \documentclass[letterpaper]{article} % DO NOT CHANGE THIS
% \usepackage{aaai2026}  % DO NOT CHANGE THIS
% \usepackage{times}  % DO NOT CHANGE THIS
% \usepackage{helvet}  % DO NOT CHANGE THIS
% \usepackage{courier}  % DO NOT CHANGE THIS
% \usepackage[hyphens]{url}  % DO NOT CHANGE THIS
% \usepackage{graphicx} % DO NOT CHANGE THIS
% \urlstyle{rm} % DO NOT CHANGE THIS
% \def\UrlFont{\rm}  % DO NOT CHANGE THIS
% \usepackage{natbib}  % DO NOT CHANGE THIS AND DO NOT ADD ANY OPTIONS TO IT
% \usepackage{caption} % DO NOT CHANGE THIS AND DO NOT ADD ANY OPTIONS TO IT
% \frenchspacing  % DO NOT CHANGE THIS
% \setlength{\pdfpagewidth}{8.5in}  % DO NOT CHANGE THIS
% \setlength{\pdfpageheight}{11in}  % DO NOT CHANGE THIS
% %
% % These are recommended to typeset algorithms but not required. See the subsubsection on algorithms. Remove them if you don't have algorithms in your paper.

% \usepackage{algorithm}
% %\usepackage{algorithmic}
% \usepackage{dsfont}
% \usepackage{algpseudocode}
% \usepackage{booktabs}
% \usepackage{multicol}
% \usepackage{multirow}
% \usepackage{enumitem}
% \usepackage{xcolor}
% \usepackage{array}
% \usepackage{makecell}
% \usepackage[most]{tcolorbox}
% \usepackage{url}
% \usepackage{adjustbox} % add this in your preamble
% \usepackage{amsmath,amssymb}
% \usepackage{amsfonts}
% %
% % These are are recommended to typeset listings but not required. See the subsubsection on listing. Remove this block if you don't have listings in your paper.
% \usepackage{newfloat}
% \usepackage{listings}
% \DeclareCaptionStyle{ruled}{labelfont=normalfont,labelsep=colon,strut=off} % DO NOT CHANGE THIS
% \lstset{%
% 	basicstyle={\footnotesize\ttfamily},% footnotesize acceptable for monospace
% 	numbers=left,numberstyle=\footnotesize,xleftmargin=2em,% show line numbers, remove this entire line if you don't want the numbers.
% 	aboveskip=0pt,belowskip=0pt,%
% 	showstringspaces=false,tabsize=2,breaklines=true}
% \floatstyle{ruled}
% \newfloat{listing}{tb}{lst}{}
% \floatname{listing}{Listing}
% %
% % Keep the \pdfinfo as shown here. There's no need
% % for you to add the /Title and /Author tags.
% \pdfinfo{
% /TemplateVersion (2026.1)
% }

% %
% % These are are recommended to typeset listings but not required. See the subsubsection on listing. Remove this block if you don't have listings in your paper.
% \usepackage{newfloat}
% \usepackage{listings}
% \DeclareCaptionStyle{ruled}{labelfont=normalfont,labelsep=colon,strut=off} % DO NOT CHANGE THIS
% \lstset{%
% 	basicstyle={\footnotesize\ttfamily},% footnotesize acceptable for monospace
% 	numbers=left,numberstyle=\footnotesize,xleftmargin=2em,% show line numbers, remove this entire line if you don't want the numbers.
% 	aboveskip=0pt,belowskip=0pt,%
% 	showstringspaces=false,tabsize=2,breaklines=true}
% \floatstyle{ruled}
% \newfloat{listing}{tb}{lst}{}
% \floatname{listing}{Listing}
% %
% % Keep the \pdfinfo as shown here. There's no need
% % for you to add the /Title and /Author tags.
% % \pdfinfo{
% % /TemplateVersion (2026.1)
% % }

% \setcounter{secnumdepth}{0} %May be changed to 1 or 2 if section numbers are desired.

% \begin{document}

% \appendix
\setcounter{secnumdepth}{1}

\renewcommand{\thetable}{A\arabic{table}}
\setcounter{table}{0}

% \onecolumn
\begin{center}
    {\LARGE \bfseries Supplementary Material}
\end{center}
\vspace{1em} 

%%%%%%%%%%%%%%%%%%%%%%%%%%%%%%%%%%%%%%%
This supplementary material provides detailed information organized as follows:
\begin{itemize}
    \item \textbf{Appendix~\ref{apd:tts_validity}}: Validity of LLM-extracted temporal annotations (adopted from our prior work (adopted from our prior work \cite{wang2025large}).
    \item \textbf{Appendix~\ref{TTS_prompt}}: LLM prompt template for extracting clinical time series from case reports.
    \item \textbf{Appendix~\ref{apd:prompts}}: Prompting strategies for zero-shot and few-shot forecasting.
    \item \textbf{Appendix~\ref{apd:encoder-forecasting}}: Technical details of encoder-based forecasting methods, including MLP-head and masking approaches.
    \item \textbf{Appendix~\ref{apd:encoder-survival}}: Survival modeling framework using language model embeddings and the hyperparameters used.
    \item \textbf{Appendix~\ref{apd:full_tables}}: Complete forecasting performance tables for all models on DeepSeek-R1 and Llama-3.3-70B annotations.
    \item \textbf{Appendix~\ref{apd:full_survival}}: Complete survival analysis performance table for all datasets.
    \item \textbf{Appendix~\ref{apd:sensitivity_analyses}}: Sensitivity analysis of the forecasting performance when using time-ordered versus text-ordered training.
    \item \textbf{Appendix~\ref{apd:masking_l33_dsr1}}: Sensitivity analysis of masking history on Llama-3.3-70B and DeepSeek-R1 annotations.
    \item \textbf{Appendix~\ref{apd:additional_insights}}: Additional insights into forecasting results.
    \item \textbf{Appendix~\ref{apd:open_source}}: Analysis of open-source models to address data leakage concerns.
    \item \textbf{Appendix~\ref{apd:beyond_pmoa}}: Generalizability beyond PubMed case reports: comparison with other clinical sources.
    \item \textbf{Appendix~\ref{apd:broader_impact}}: Broader societal impact and ethical considerations.
\end{itemize}

%%%%%%%%%%%%%%%%%%%%%%%%%%%%%%%%%%%%%%%%%%%%%%%%
%%%%%%%%%Appendix A: T2S2 extraction pipeline
%%%%%%%%%%%%%%%%%%%%%%%%%%%%%%%%%%%%%%%%%%%%%%%
\section{Validity of Temporal Annotations}
\label{apd:tts_validity}

%\subsection{Evaluation Metrics}

We used \texttt{DeepSeek-R1-UD-IQ1} and \texttt{Llama-3.3-70B-Instruct} models to generate clinical textual time series from
the 2319 sepsis case reports. We refer to this as Textual
Time Series for Sepsis (\textit{T2S2}). We assessed the quality of \textit{T2S2} using three key metrics: event match rate, temporal concordance (c-index), and log-time discrepancy on two external validation sets: \textit{sepsis-10} (n = 10) and \textit{sepsis-100} (n = 90).(adopted from our prior work \cite{wang2025large}).

\subsection*{Text Event Match}
%\textbf{Text Event Match.} 
To align predicted clinical findings with reference annotations, we use a recursive best-match procedure adapted from \citep{wang2025large}. At each step, we select the reference-prediction pair with the minimal string-based distance; ties are broken by text order (\emph{i.e.}, the order of events in the annotation files). Matched items are removed and the process recurses on the remainder, preventing many-to-one mappings. We evaluated Levenshtein distance, cosine similarity of \textit{BERT} (\textit{bert-base-uncased}) embeddings, and sentence-embedding cosine similarity from \textit{PubMedBERT} (\textit{S-PubMedBert-MS-MARCO}). Prior validation showed \textit{PubMedBERT-based} cosine similarity with a threshold of 0.1 performed best, and we used this metric for evaluation. Pseudocode is provided in Algorithm \ref{alg:recursive_match}.

%\input{Scientific_data/Supplementary/recursive_best_match_algo}

\subsection*{Temporal Ordering}
%\textbf{Temporal Ordering.} 
% (Optional) Requires \usepackage{dsfont} for \mathds{1}
Temporal ordering accuracy was quantified with the concordance index (c-index), the probability that a randomly chosen pair of events is correctly ordered. For reference times $\{t^r_i\}$ and predicted times $\{t^p_i\}$,
\[
\text{c-index}=\frac{1}{N}\sum_{i<j}\mathds{1}\{(t^r_i-t^r_j)(t^p_i-t^p_j)>0\},
\]
where $N$ is the number of comparable pairs (event indices $i,j$ with $t^r_i\neq t^r_j$ and $t^p_i\neq t^p_j$), and $\mathds{1}\{\cdot\}$ is the indicator function. Higher values indicate better agreement with the reference timeline.
 


\begin{algorithm}[tbp]
\caption{Recursive Best Match}
\label{alg:recursive_match}
\begin{algorithmic}[1]
\Require Two lists: \texttt{ref} (reference events) and \texttt{pred} (predicted events)
\Ensure List of best-matching event pairs
\Function{MatchEvents}{\texttt{ref}, \texttt{pred}}
    \If{\texttt{ref} is empty \textbf{or} \texttt{pred} is empty}
        \State \Return $[]$
    \EndIf
    \State $\text{min\_distance} \gets \infty$
    \State $\text{best\_pair} \gets \text{None}$
    \ForAll{$r \in \texttt{ref}$}
        \ForAll{$p \in \texttt{pred}$}
            \State $d \gets \text{ComputeDistance}(r, p)$
            \If{$d < \text{min\_distance}$}
                \State $\text{min\_distance} \gets d$
                \State $\text{best\_pair} \gets (r, p)$
            \ElsIf{$d = \text{min\_distance}$}
                \State Get indices of $r$, $p$, and current $\text{best\_pair}$
                \If{index of $r <$ best $r$ index}
                    \State $\text{best\_pair} \gets (r, p)$
                \ElsIf{index of $r = $ best $r$ index \textbf{and} index of $p <$ best $p$ index}
                    \State $\text{best\_pair} \gets (r, p)$
                \EndIf
            \EndIf
        \EndFor
    \EndFor
    \State Remove $\text{best\_pair.r}$ from \texttt{ref}
    \State Remove $\text{best\_pair.p}$ from \texttt{pred}
    \State $\text{result} \gets [\text{best\_pair}] + \text{MatchEvents}(\texttt{ref}, \texttt{pred})$
    \State \Return $\text{result}$
\EndFunction
\end{algorithmic}
\end{algorithm}
%\vspace{-5pt}


\subsection*{Time Discrepancy Analysis}
%\textbf{Time Discrepancy Analysis.} 
 Temporal prediction accuracy was assessed via discrepancies between predicted and reference timestamps. We computed the log-time discrepancy as $\log(1+\Delta t)$, where $\Delta t$ is the absolute time difference (in hours), and summarized errors with the Area Under the Log-Time Cumulative Distribution Function (AULTC), which reflects how rapidly the error CDF rises—larger AULTC indicates better temporal alignment. 

% One-column AAAI-friendly table (stacked datasets)
\begin{table}[thbp]
\centering
% \scriptsize
\begin{adjustbox}{max width=\columnwidth}
\setlength{\tabcolsep}{4.5pt} % tighten to fit one column
\renewcommand{\arraystretch}{1.1}

\begin{tabular}{lccc}
\hline
\multicolumn{4}{c}{\textbf{\textit{sepsis-10}}} \\
\hline
\textbf{Model} & \textbf{Event Match Rate} & \textbf{Concordance} & \textbf{AULTC} \\
\texttt{DeepSeek-R1}      & 0.79 & 0.876 & 0.772 \\
\texttt{Llama-3.3-70B}    & 0.75 & 0.932 & 0.759 \\
\hline
\multicolumn{4}{c}{\textbf{\textit{sepsis-100}}} \\
\hline
\textbf{Model} & \textbf{Event Match Rate} & \textbf{Concordance} & \textbf{AULTC} \\
\texttt{DeepSeek-R1}      & 0.91 & 0.875 & 0.801 \\
\texttt{Llama-3.3-70B}    & 0.78 & 0.894 & 0.791 \\
\hline
\end{tabular}
\end{adjustbox}
\caption{Comparison of event match rate (Event), median concordance (c), median absolute error (MAE), and area under the log-time curve (AULTC) across \texttt{DeepSeek-R1} and \texttt{Llama-3.3-70B} for \textit{sepsis-10}, and \textit{sepsis-100} datasets.}
\label{tab:sepsis_onecol}
\end{table}


Formally, we compute the AULTC as the area under $F(x)$ from $x = 0$ to $x = \log(1 + S_{\text{max}})$, normalized by $\log(1 + S_{\text{max}})$:

\begin{equation*}
\label{eq:aultc}
\begin{aligned}
\text{AULTC}
&= \frac{1}{\log(1+S_{\max})}
   \sum_{i=1}^{k}\big(x_{(i)}-x_{(i-1)}\big)\frac{i}{k} \\
&\quad + \frac{\log(1+S_{\max})-x_{(k)}}{\log(1+S_{\max})},
\end{aligned}
\end{equation*}

where $x_{(0)} = 0$. With this definition, AULTC = 1 indicates that discrepancies are zero (perfect recovery), resulting in maximum area $\log(1 + S_{\text{max}})$, and AULTC = 0 indicates that all discrepancies exceed $S_{\text{max}}$, yielding zero area. We also stratified results by reference time intervals (hour, day, week, year) to examine accuracy as a function of time from case presentation.

\subsection*{Annotation Quality Results}
On the \textit{sepsis-10} subset, \texttt{DeepSeek-R1} achieved an event match rate of 0.79, c-index of 0.876, and AULTC of 0.772, while \texttt{Llama-3.3-70B} obtained 0.75, 0.932, and 0.759, respectively. Thus, \texttt{Llama} provided stronger pairwise ordering (higher c-index) but with fewer matched events and a slightly weaker log-time error profile (lower AULTC) than \texttt{DeepSeek}. On \textit{sepsis-100}, \texttt{DeepSeek-R1} reached 0.91 (event match rate), 0.875 (c-index), and 0.801 (AULTC), whereas \texttt{Llama-3.3-70B} reached 0.78, 0.894, and 0.791. Here, \texttt{DeepSeek} matched substantially more events and achieved better AULTC, while \texttt{Llama} again showed a small c-index advantage. Overall, \texttt{DeepSeek} tends to optimize event identification and log-time alignment, whereas \texttt{Llama} slightly favors pairwise temporal ordering. 

Overall, these metrics together indicate strong alignment of the generated textual time series with expert annotations, sufficient to support forecasting benchmarks. 


%%%%%%%%%%%%%%%%%%%%%%%%%%%%%%%%%%%%%%%%%%%%%%%%
%%%%%%%%%Appendix B: LLM prompt
%%%%%%%%%%%%%%%%%%%%%%%%%%%%%%%%%%%%%%%%%%%%%%%
\onecolumn

{\refstepcounter{section} \label{TTS_prompt}
\Large\bfseries\thesection\quad LLM Prompt to Generate Clinical Textual Time Series from PMOA Case Reports
}

\vspace{1em}

% \begin{center}
\begin{minipage}[!htbp]{0.99\textwidth}
\begin{tcolorbox}[colback=gray!5, colframe=gray!100, title=Prompt, boxrule=2pt]

{You are a physician. 

Extract the clinical events and the related time stamp from the case report. The admission event has timestamp 0. If the event is not available, we treat the event, e.g. current main clinical diagnosis or treatment with timestamp 0. The events that happened before event with 0 timestamp have negative time, the ones after the event with 0 timestamp have positive time. The timestamps are in hours. The unit will be omitted when output the result. If there is no temporal information of the event, please use your knowledge and events with temporal expression before and after the events to provide an approximation. We want to predict the future events given the events happened in history. For example, here is the case report.}\\

\noindent
\textcolor{teal}{\texttt{An 18-year-old male was admitted to the hospital with a 3-day history of fever and rash. Four weeks ago, he was diagnosed with acne and received the treatment with minocycline, 100 mg daily, for 3 weeks. With increased WBC count, eosinophilia, and systemic involvement, this patient was diagnosed with DRESS syndrome. The fever and rash persisted through admission, and diffuse erythematous or maculopapular eruption with pruritus was present. One day later the patient was discharged.}}\\

\noindent
{Let's find the locations of event in the case report, it shows that four weeks ago of fever and rash, four weeks ago, he was diagnosed with acne and receive treatment. So the event of fever and rash happened four weeks ago, 672 hours, it is before admitted to the hospital, so the time stamp is -672. diffuse erythematous or maculopapular eruption with pruritus was documented on the admission exam, so the timestamp is 0 hours, since it happens right at admission. DRESS syndrome has no specific time, but it should happen soon after admission to the hospital, so we use our clinical judgment to give the diagnosis of DRESS syndrome the timestamp 0. then the output should look like:}\\
\vspace{2mm}
\noindent 

\texttt{18 years old | 0 \\
male | 0 \\
admitted to the hospital | 0 \\
fever | -72 \\
rash | -72 \\
acne | -672 \\
minocycline | -672 \\
increased WBC count | 0 \\
eosinophilia | 0 \\
systemic involvement | 0 \\
diffuse erythematous or maculopapular eruption | 0 \\
pruritus | 0 \\
DRESS syndrome | 0 \\
fever persisted | 0 \\
rash persisted | 0 \\
discharged | 24}\\

\noindent {Separate conjunctive phrases into its component events and assign them the same timestamp (for example, separation of ‘fever and rash’ into 2 events: ‘fever’ and ‘rash’). If the event has duration, assign the event time as the start of the time interval. Attempt to use the text span without modifications except ‘history of’ where applicable. Include all patient events, even if they appear in the discussion; do not omit any events; include termination/discontinuation events; include the pertinent negative findings, like ‘no shortness of breath’ and ‘denies chest pain’. Show the events and timestamps in rows, each row has two columns: one column for the event, the other column for the timestamp. The time is a numeric value in hour unit. The two columns are separated by a pipe $\mid$ as a bar-separated file. Skip the title of the table. Reply with the table only.}

\end{tcolorbox}
\end{minipage}
% \end{center}

%%%%%%%%%%%%%%%%%%%%%%%%%%%%%%%%%%%%%%%%%%%%%%%%
%%%%%%%%%Appendix C: Prompting strategies
%%%%%%%%%%%%%%%%%%%%%%%%%%%%%%%%%%%%%%%%%%%%%%%
\twocolumn

\onecolumn

{\refstepcounter{section}
\noindent\Large\bfseries\thesection\quad Prompting Strategies for Forecasting Tasks
\label{apd:prompts}}

\subsection{Zero-Shot}

\begin{tcolorbox}[colback=gray!5, colframe=gray!100, title=Prompt, boxrule=2pt]

{You are an expert physician.}

{Reply to the prompt with structured predictions in a k-item, bar-separated row.

For example, if there are k=3 events (A, B, C) and only B occurs in the time window, then:}
{0 $\mid$ 1 $\mid$  0
would be the correct response. \\ 

Reply with the answer only. Do NOT provide any other text. Do NOT explain.}

\end{tcolorbox}

\subsection{Zero-Shot (with degree of confidence)}

\begin{tcolorbox}[colback=gray!5, colframe=gray!100, title=Prompt, boxrule=2pt]
{You are an expert physician.}

{Reply to the prompt with structured predictions in a k-item, bar-separated row.
Assess whether the event happens in the time window using a scale from 1 to 5, where 5 means it definitely happened and 1 means it did not happen.

For example, if there are k=3 events (A, B, C) and only B occurs in the time window, then:
$1 | 5 | 1$
would be the correct response. \\

Reply with the answer only. Do NOT provide any other text. Do NOT explain.}
\end{tcolorbox}

\subsection{Few-Shot}
\begin{tcolorbox}[colback=gray!5, colframe=gray!100, title=Prompt, boxrule=2pt]

{You are an expert physician.}

{Reply to the prompt with predictions in a k-item, bar-separated row to answer if the k events happen within a given forecast time window. \\

Example 1. If there are k=3 events (A, B, C) and only B occurs in the requested time window, then:
0 $\mid$ 1 $\mid$ 0
would be the correct response.}\\

{Example 2. Suppose the time series is:}

\noindent \texttt{\textcolor{teal}
{admitted | 0 \\
history of diabetes | 0 \\
elevated troponin | 5 \\
--- \textbf{PREDICTION WINDOW} --- \\
cardiac catheterization | 6 \\
admission to icu | 24 \\ 
transfer to floor | 72\\ 
discharge | 144\\ 
office visit | 1440}
}

{Then, if the prediction is at 5 hours, the time preceding ``--- PREDICTION WINDOW ---'', the events admitted, history of diabetes, and elevated troponin have occurred.  Supposing k=4 and a forecast window of 24 hours, then the task it to determine which of cardiac catheterization, admission to icu, transfer to floor, and discharge happen. So the correct response for the example is:
1 $\mid$ 1 $\mid$ 0 $\mid$ 0
because cardiac catherization (6) and admission to icu (24) happen before 5 + 24 = 29.}\\

{Reply with the answer only. Do NOT provide any other text. Do NOT explain.}

\end{tcolorbox}

\subsection{Few-Shot Ordinal}

\begin{tcolorbox}[colback=gray!5, colframe=gray!100, title=Prompt, boxrule=2pt]

{You are an expert physician.}

{Reply to the prompt with predictions in a k-item, bar-separated row to answer if the k events happen within a given forecast time window, using a scale from 1 to 5, where 5 means it definitely happened and 1 means it did not happen.}\\

{Example 1. If there are k=3 events (A, B, C) and only B occurs in the time window, then:
1 $\mid$ 5 $\mid$ 1
would be the correct response.}\\

{Example 2. Suppose the time series is:}\\
\noindent \texttt{\textcolor{teal}{
admitted | 0\\
history of diabetes | 0\\
elevated troponin | 5\\
--- \textbf{PREDICTION WINDOW} ---\\
cardiac catheterization | 6\\
admission to icu | 24\\
transfer to floor | 72\\
discharge | 144\\
office visit | 1440}}

{Then, if the prediction is at 5 hours, the time preceding ``--- PREDICTION WINDOW ---'', the events admitted, history of diabetes, and elevated troponin have occurred.  Supposing k=4 and a forecast window of 24 hours, then the task it to determine which of cardiac catheterization, admission to icu, transfer to floor and discharge happen. So the an excellent response for the example is:
5 $\mid$ 4 $\mid$ 1 $\mid$ 1
because cardiac catherization (6) and admission to icu (24) happen before 5 + 24 = 29.}\\

{Reply with the answer only. Do NOT provide any other text. Do NOT explain.}

\end{tcolorbox}

\twocolumn

%%%%%%%%%%%%%%%%%%%%%%%%%%%%%%%%%%%%%%%%%%%%%%%%
%%%%%%%%%Appendix D: Textual Time Series Forecasting Framework of Encoder Models
%%%%%%%%%%%%%%%%%%%%%%%%%%%%%%%%%%%%%%%%%%%%%%%

\clearpage
\section{Textual Time Series Forecasting Framework of Encoder Models}
\label{apd:encoder-forecasting}

We develop a general framework for predictive modeling over textual time series of the encoder-only models, designed to accommodate the irregular, event-based structure common to clinical case reports. 
As highlighted Section~\ref{sec:tts-annotation}, each data instance consists of a sequence of clinical events, where each event is a free-text string describing a clinical finding, intervention, or diagnosis, paired with a timestamp in hours relative to the time of admission. 
The format resembles a two-column table, where the first column contains textual clinical events and the second column contains integer-valued timestamps.

\paragraph{Input Preprocessing and Timeseries Construction.}
From each case report, we construct multiple training examples for sequence forecasting. 
For a fixed forecast horizon parameter $K$, we generate overlapping windows of historical events and forecast targets. 
Specifically, for each case report with $L$ rows, we construct a sequence of up to $L-K$ forecasting examples. 
Each example is structured as follows:

\begin{itemize}
    \item A \textbf{history segment} $\mathcal{H}_t = \{(e_i, t_i)\}_{i=1}^{n}$, consisting of clinical events $e_i$ occurring at timestamps $t_i \leq t$.
    \item A \textbf{forecast target} $\mathcal{F}_t = \{(f_j, \tau_j)\}_{j=1}^{K}$, consisting of the next $K$ events after $t$.
\end{itemize}

To preserve the temporal coherence of events, we treat all events with the same timestamp as an atomic unit. 
That is, if multiple events share the same timestamp, they are included together in the history or forecast block, and the next forecasting window is only advanced after the full timestamp block has been consumed.

\paragraph{Model Input Representation.}
Each training example is serialized into a transformer-compatible input by linearizing the historical events and forecast events. 
Importantly, we preserve both temporal and textual content by encoding each event in the form ``\texttt{t:e}'' where \texttt{t} is the timestamp and \texttt{e} is the event string.
The input to the language model takes the form:

\begin{equation}
\texttt{[CLS]} \; t_1\texttt{:}e_1 \; \texttt{[SEP]} \; t_2\texttt{:}e_2 \; \texttt{[SEP]} \; \dots \; \texttt{[SEP]} \; t_n\texttt{:}e_n
\end{equation}

This is followed by additional target-specific components, as described below.

\paragraph{Language Models.}
To test the generality of our framework, we experiment with a variety of pre-trained encoder-only language models, including:
\begin{itemize}
    \item \textit{bert-base-uncased} (BERT; Devlin et al., 2018)
    \item \textit{roberta-base} (Liu et al., 2019)
    \item \textit{microsoft/deberta-v3-small} (He et al., 2021)
    \item \textit{answerdotai/ModernBERT-base}, \textit{answerdotai/ModernBERT-large}
    \item \textit{thomas-sounack/BioClinical-ModernBERT-base}, \textit{thomas-sounack/BioClinical-ModernBERT-large} 
\end{itemize}
Each model is used with its corresponding tokenizer. 
For models with SentencePiece tokenization (e.g., \textit{DeBERTa}), we use the appropriate fast tokenizer such as \textit{DebertaV2Tokenizer} to ensure compatibility. 
All inputs are truncated from the left (i.e., keeping the most recent tokens) if they exceed the model's context window (typically 512 tokens).

\paragraph{Prediction Tasks.}
We define two primary prediction tasks over $\mathcal{F}_t$, each formulated as a binary classification problem.

\subparagraph{1. Event Ordering (Concordance) Task.}
The goal in this task is to predict the correct temporal ordering between future events. 
For each example, we generate all \emph{comparable pairs} of events in $\mathcal{F}_t$, defined as pairs $(f_i, f_j)$ where $\tau_i \neq \tau_j$. 
For each pair, we form a binary classification example as follows:

\begin{itemize}
    \item \textbf{Input:} The serialized history segment followed by:
    \begin{equation}
        \texttt{[SEP]} \; \tau_i\texttt{:}f_i \; \texttt{[SEP]} \; \tau_j\texttt{:}f_j \; \texttt{[SEP]}
    \end{equation}
    \item \textbf{Label:} $y = 1$ if $\tau_i < \tau_j$, else $y = 0$
\end{itemize}

The model output is taken from the \texttt{[CLS]} token, passed through a linear classification head, and softmaxed to produce a binary probability. 
We optimize cross-entropy loss across all pairwise comparisons. 
At evaluation time, we report the concordance index (c-index), which measures the proportion of correctly predicted pairs among all comparable pairs.

\subparagraph{2. Time-Window Classification Task.}
This task evaluates whether each forecast event occurs within a fixed prediction window $H$ hours of the last historical event. 
For each $f_j \in \mathcal{F}_t$, we generate an input:

\begin{itemize}
    \item \textbf{Input:} The serialized history segment followed by:
    \begin{equation}
        \texttt{[SEP]} \; \tau_j\texttt{:}f_j \; \texttt{[SEP]}
    \end{equation}
    \item \textbf{Label:} $y = 1$ if $\tau_j - t_n \leq H$, where $t_n$ is the timestamp of the last historical event.
\end{itemize}

This is again treated as a binary classification task over the \texttt{[CLS]} embedding, using a linear layer and softmax. 
We report macro-averaged F1 scores over all prediction examples.

\paragraph{Masked Language Modeling Approaches.}
In addition to the MLP-head approach described above, we implement masked language modeling (MLM) variants that leverage the pretrained masked token prediction capabilities of encoder models. These approaches reformulate our forecasting tasks as cloze-style problems:

\subparagraph{1. Masked Event Occurrence Prediction.}
For the time-window classification task, we construct prompts that include a \texttt{[MASK]} token whose prediction indicates whether an event occurs within the specified horizon:

\begin{itemize}
    \item \textbf{Input:} History formatted as: \texttt{[CLS]} $h_1$ \texttt{[SEP]} $h_2$ \texttt{[SEP]} ... \texttt{[SEP]} Will "$e_j$" happen within $H$ hours? \texttt{[MASK]} \texttt{[SEP]}
    \item \textbf{Target:} The model predicts either "\texttt{yes}" or "\texttt{no}" at the \texttt{[MASK]} position
\end{itemize}

We constrain the output vocabulary to only these two tokens and compute cross-entropy loss over their logits. For fine-tuned variants, we train the entire model to optimize this masked prediction task. For zero-shot variants, we directly use the pretrained model's predictions without additional training.

\subparagraph{2. Masked Temporal Ordering.}
For the concordance task, we predict relational tokens that indicate temporal ordering:

\begin{itemize}
    \item \textbf{Input:} History followed by: \texttt{[SEP]} $e_1$ \texttt{[MASK]} $e_2$ \texttt{[SEP]}
    \item \textbf{Target:} The model predicts "\texttt{before}" if $e_1$ occurs before $e_2$, else "\texttt{after}"
\end{itemize}

Again, we restrict predictions to these two tokens. This formulation allows the model to leverage its pretrained understanding of temporal relationships encoded during MLM pretraining. 

Note that in our dataset, we randomize the relative ordering of $e_1$ and $e_2$ each time to avoid introducing positional biases that could lead the model to rely on superficial cues rather than genuine temporal reasoning.

\subparagraph{Implementation Details.}
For both tasks, we extract the logits at the \texttt{[MASK]} position and apply a linear projection to the target vocabulary size (2 tokens). The key advantage of this approach is that it can leverage pretrained masked language modeling capabilities, potentially requiring less task-specific training. We evaluate both fine-tuned versions (which update all model parameters) and zero-shot versions (which use frozen pretrained weights).

\paragraph{Training and Evaluation.}
Each task is trained separately using AdamW with linear learning rate decay and early stopping based on validation loss. 
We save the model checkpoint with the best validation performance. 
During inference, we reload this checkpoint and evaluate on a held-out test set (which is obtained from textual time series of unseen case reports).

For both tasks, we extract softmax probabilities rather than hard labels to allow for threshold calibration and ROC-based analyses. 
All training is conducted using PyTorch and HuggingFace Transformers.
%%%%%%%%%%%%%%%%%%%%%%%%%%%%%%%%%%%%%%%%%%%%%%%%
%%%%%%%%%Appendix E: Survival Modeling with Language Model Embeddings
%%%%%%%%%%%%%%%%%%%%%%%%%%%%%%%%%%%%%%%%%%%%%%%
\clearpage
\section{Survival Modeling with Language Model Embeddings}
\label{apd:encoder-survival}

\begin{table}[b]
\centering
\begin{adjustbox}{max width=\columnwidth}
\begin{tabular}{ll}
\toprule
\textbf{Model} & \textbf{Hyperparameters} \\
\midrule
RSF & \texttt{n\_estimators} $\in$ \{20,50,100,200,500,1000\} \\
    & \texttt{min\_samples\_split} $\in$ \{2,10,20,50,100\} \\
    & \texttt{min\_samples\_leaf} $\in$ \{1,5,10,20\} \\
DeepSurv/DeepHit & \texttt{num\_nodes} $\in$ \{32,64,128,256,512,768,1024,1200,\\
& 1536,1800,2048,3072,4096\} \\
    & \texttt{dropout} $\in$ \{0.1,0.2,0.4,0.5\} \\
\bottomrule
\end{tabular}
\end{adjustbox}
\caption{Hyperparameter grid search ranges for our survival models.}
\label{tab:hyperparams}
\end{table}

To evaluate the prognostic value of textual information encoded in large language models (LLMs), we adopt a two-stage framework: (1) extraction of fixed-dimensional sequence embeddings from various pre-trained LLMs, and (2) downstream survival modeling using these embeddings as covariates.

In the first stage, we process each textual time series (which is converted to a textual context of \emph{"(time) clinical event [SEP] ... [SEP] (time) clinical event [SEP]"}) using a suite of LLMs to obtain dense vector representations that summarize the content of the input sequence. 
For models belonging to the encoder family, including \textit{bert-base-uncased}, \textit{roberta-base}, \textit{deberta-v3-small}, \textit{ModernBERT-base}, and \textit{ModernBERT-large}, we extract the final hidden state corresponding to the \textit{[CLS]} token, which is conventionally used to represent the entire sequence in classification tasks. 
This token-specific embedding serves as a compact, sequence-level representation.

For decoder-based models that do not utilize a dedicated \texttt{[CLS]} token, including \textit{DeepSeek-R1-Distill-Llama-70B}, \textit{Llama-3.3-70B-Instruct}, \textit{DeepSeek-R1-Distill-Llama-8B}, and \textit{Llama-3.1-8B-Instruct}, we compute the mean-pooled embedding over the last hidden states of all non-padding tokens. 
This pooling strategy yields a fixed-length vector that captures the overall semantic content of the input while mitigating the impact of padding artifacts.

In the second stage, these embedding vectors are used as input covariates to three survival models, each designed to capture time-to-event dynamics in different ways:

\begin{itemize}
    \item \textbf{Random Survival Forest (RSF):} 
    A nonparametric ensemble method based on decision trees, capable of handling complex interactions and high-dimensional inputs. 
    The model is tuned over the number of trees (\texttt{n\_estimators}), the minimum number of samples required to split an internal node (\texttt{min\_samples\_split}), and the minimum number of samples at a leaf node (\texttt{min\_samples\_leaf}); see Table~\ref{tab:hyperparams} for search ranges.

    \item \textbf{DeepSurv:} 
    A neural-network generalization of the Cox proportional hazards model. 
    The architecture comprises fully connected layers with varying widths (\texttt{num\_nodes}) and dropout rates; see Table~\ref{tab:hyperparams} for search ranges. The number of training epochs is fixed.

    \item \textbf{DeepHit:} 
    A multi-task neural network that jointly models the discrete hazard and survival functions. 
    As with DeepSurv, we perform a grid search over hidden layer sizes and dropout rates while keeping the number of epochs fixed; see Table~\ref{tab:hyperparams}.
\end{itemize}

Hyperparameter tuning for each survival model is performed using a validation split from the training set. 
For each combination of language model and survival model, we conduct grid search over the respective hyperparameter space, selecting the configuration that maximizes the time-dependent concordance index \citep{antolini2005time} on the validation set. 
Final evaluation is conducted on held-out test sets (\textit{T2S2\_test}, \textit{sepsis-10}, \textit{sepsis-100}), and results are aggregated across survival metrics to assess model performance.

This framework allows for a systematic evaluation of how different pre-trained LLMs, and their associated embedding strategies, contribute to downstream survival prediction tasks.
%%%%%%%%%%%%%%%%%%%%%%%%%%%%%%%%%%%%%%%%%%%%%%%%
%%%%%%%%%Appendix F: Forecasting tables
%%%%%%%%%%%%%%%%%%%%%%%%%%%%%%%%%%%%%%%%%%%%%%%
\clearpage
\onecolumn
% \section{Forecasting performance on \textit{DeepSeek-R1} and \textit{Llama-3.3-70B} textual time series annotations}
% \label{apd:full_tables}

{\refstepcounter{section}
\noindent\Large\bfseries\thesection\quad Forecasting Performance on \texttt{DeepSeek-R1} and \texttt{Llama-3.3} Textual Time Series Annotations
\label{apd:full_tables}}

\vspace{0em}

Tables~\ref{tab:DS_forecasting} and~\ref{tab:L33_forecasting} present the forecasting performance of our models on textual time series annotations generated by \texttt{DeepSeek-R1} and \texttt{Llama-3.3-70B}, respectively. We report standard forecasting metrics across multiple evaluation settings to provide a comprehensive view of model performance. 

\input{AAAI_2026/Tables/DS_time_ordered}
\input{AAAI_2026/Tables/L33_time_ordered}


%%%%%%%%%%%%%%%%%%%%%%%%%%%%%%%%%%%%%%%%%%%%%%%%%%%%%%%%%%%%%%%%
\clearpage
% \twocolumn
{\refstepcounter{section}
\noindent\Large\bfseries\thesection\quad Survival Performance on \texttt{DeepSeek-R1} Annotations for \textit{T2S2}, \textit{sepsis-10}, and \textit{sepsis-100}
\label{apd:full_survival}}
% \section{Survival performance on \textit{DeepSeek-R1} annotations for T2S2, sepsis-10, and \textit{sepsis-100}}
% \label{apd:full_survival}
\vspace{1em}

In this section, we report the full concordance analyses for the \textit{sepsis-10} and \textit{sepsis-100} annotation sets in addition to \textit{T2S2} test set, presented in Table~\ref{tab:concordance_wide}. These results are the extended version of the \textit{T2S2}-only table (Table~\ref{tab:concordance_t2s2_only}) included in the main text.

\input{AAAI_2026/Tables/survival_table_full}


%%%%%%%%%%%%%%%%%%%%%%%%%%%%%%%%%%%%%%%%%%%%%%%%%%%%%%%%%%%%%%%%% Appendix G: Masking history
%%%%%%%%%%%%%%%%%%%%%%%%%%%%%%%%%%%%%%%%%%%%%%%%%%%%%%%%%%%%%%%%
% \clearpage
% \section{Masking history: \textit{Llama-3.3-70B} textual time series annotations}
% \label{apd:masking_l33}
\vspace{2em}

{\refstepcounter{section}
\noindent\Large\bfseries\thesection\quad Sensitivity Analysis of Time-Ordered versus Text-Ordered Training
\label{apd:sensitivity_analyses}}

\vspace{1em}
% \section{Sensitivity analysis}
% \vspace{-0.2em}
% \paragraph{Time-Ordered versus Text-Ordered Training.}

Table \ref{tab:dsr1q_l33_comparison} shows time-ordering generally improves concordance, particularly for \texttt{DeepSeek-R1} annotations, while F1 scores show mixed results. Text-ordering sometimes yields better external dataset performance with \texttt{Llama-3.3-70B}, suggesting that preserving the narrative sequence can support accurate classification even when it does not follow temporal order.  These results highlight a trade-off between temporal alignment and narrative coherence for different tasks.

\input{AAAI_2026/Tables/tab_textord_vs_timeord}

%%%%%%%%%%%%%%%%%%%%%%%%%%%%%%%%%%%%%%%%%%%%%%%%
%%%%%%%%%%%%%%%%%%%%%%%%%%%%%%%%%%%%%%%%%%%%%%%%
\clearpage
{\refstepcounter{section}
\noindent\Large\bfseries\thesection\quad Masking History: \texttt{Llama-3.3-70B} Textual Time Series Annotations
\label{apd:masking_l33_dsr1}}

\vspace{1em}

% Table~\ref{tab:masking_effects_llama} examines the impact of randomly masking portions of the input history on model performance when using \textit{Llama-3.3-70B}-generated annotations. 
% We evaluate performance across three datasets (T2S2-L33, Sepsis-10, and \textit{sepsis-100}) using F1 scores for prediction of next \(k\) events happening within the next day and concordance index (c-index) for correctly ordering the next \(k\) events. 
% As the proportion of masked history increases, we observe a general decline in F1 performance, particularly at higher masking rates (e.g., 90\%), suggesting that access to historical context is important for accurate forecasting. 
% Interestingly, the concordance index remains relatively stable, indicating the model's robustness in ranking the next \(k\) events with partial input histories.


% \paragraph{Impact of timestep dropout on F1 and concordance.}
Tables~\ref{tab:masking_effects_DS} and~\ref{tab:masking_effects_llama} report the effect of timestep dropout on classification and ranking performance. 
We evaluate performance across three datasets (\textit{T2S2}-\texttt{L3.3}, \textit{sepsis-10}, and \textit{sepsis-100}) using F1 scores for prediction of next \(k\) events happening within the next day and concordance index (c-index) for correctly ordering the next \(k\) events. 
As the proportion of masked history increases, we observe a general decline in F1 performance, particularly at higher masking rates (e.g., 90\%), suggesting that access to historical context is important for accurate forecasting. 
Interestingly, the concordance index remains relatively stable, indicating the model's robustness in ranking the next \(k\) events with partial input histories.

More specifically, we can see that F1 scores decrease monotonically with increasing TDR, with relatively mild degradation up to 30--60\% dropout and sharper declines thereafter. For example, in the \texttt{DeepSeek-R1} annotations (Table~\ref{tab:masking_effects_DS}), \textit{T2S2}/\textit{sepsis-10}/\textit{sepsis-100} drop from 0.576/0.428/0.607 at TDR = 0\% to 0.423/0.200/0.510 at TDR = 90\%, with \textit{sepsis-10} showing the largest relative decline. In contrast, concordance (c-index) remains stable across dropout levels (0.661--0.678), indicating that ranking-based evaluations are less sensitive to history truncation. These trends are consistent for \texttt{Llama-3.3} annotations (Table~\ref{tab:masking_effects_llama}), suggesting that ranking objectives capture temporal relationships that are more robust to missing historical information than binary classification accuracy of predicting whether an event happens within a time window.

\begin{table}[ht]
    \centering
    % \resizebox{\linewidth}{!}{
    \begin{tabular}{lcccccc}
        \toprule
        \textbf{Timestep drop rate} & F1(1d) & c-index & F1(1d)-10 & c10 & F1(1d)-100 & c100 \\
        \hline
        0\%  & 0.576 & 0.672 & 0.428 & 0.576 & 0.607 & 0.613 \\
        15\% & 0.572 & 0.667 & 0.415 & 0.564 & 0.604 & 0.616 \\
        30\% & 0.548 & 0.672 & 0.357 & 0.542 & 0.591 & 0.619 \\
        60\% & 0.542 & 0.678 & 0.310 & 0.578 & 0.588 & 0.616 \\
        90\% & 0.423 & 0.661 & 0.200 & 0.561 & 0.510 & 0.614 \\
        \bottomrule
    \end{tabular}
    % }
    \caption{Effect of Randomly Masking History (\texttt{DeepSeek-R1} annotations) on F1 and Concordance. \textbf{Abbreviations}: F1(1d)-10/100: F1 (1 day) for \textit{sepsis-10} and \textit{sepsis-100} respectively. c10/100: c-index for \textit{sepsis-10} and \textit{sepsis-100} respectively.}
    \label{tab:masking_effects_DS}
    % \vspace{-1.5em}
\end{table}

\begin{table*}[!hbpt]
    \centering
    \resizebox{0.75\linewidth}{!}{
    \begin{tabular}{lcccccc}
        \toprule
        \textbf{Timestep drop rate} & F1(1d) & c-index & F1(1d)-10 & c10 & F1(1d)-100 & c100 \\
        \hline
        0\%  & 0.645 & 0.632 & 0.431 & 0.595 & 0.626 & 0.612 \\
        15\% & 0.628 & 0.635 & 0.422 & 0.592 & 0.624 & 0.610 \\
        30\% & 0.634 & 0.634 & 0.423 & 0.586 & 0.621 & 0.598 \\
        60\% & 0.604 & 0.634 & 0.391 & 0.577 & 0.605 & 0.613 \\
        90\% & 0.531 & 0.622 & 0.328 & 0.546 & 0.534 & 0.607 \\
        \bottomrule
    \end{tabular}
    }
    \caption{Effect of Randomly Masking History (\texttt{Llama-3.3-70B} annotations) on F1 and Concordance. \textbf{Abbreviations}: F1(1d)-10/100: F1 (1 day) for \textit{sepsis-10} and \textit{sepsis-100} respectively. c10/100: c-index for \textit{sepsis-10} and \textit{sepsis-100} respectively.}
    \label{tab:masking_effects_llama}
\end{table*}

%%%%%%%%%%%%%%%%%%%%%%%%%%%%%%%%%%%%%%%%%%%%%%%%
%%%%%%%%%Appendix H:  
%%%%%%%%%%%%%%%%%%%%%%%%%%%%%%%%%%%%%%%%%%%%%%%
\clearpage
\twocolumn
\section{Additional Insights into Forecasting Results}
\label{apd:additional_insights}

In this section, we provide some additional insights into the forecasting results. Specifically, we summarize two observations that clarify model behavior and apparent anomalies in the results.

\subsection{Why decoders fare better on survival but not short-horizon forecasting}
 Our findings align with adjacent literature on representation usage over different time horizons. \emph{Encoder} models, typically pooled via a \texttt{[CLS]} token, excel on short-horizon tasks because bidirectional attention aggregates local context effectively \citep{devlin2019bert,bai2018empirical}. In contrast, \emph{decoder} models trained autoregressively pass information forward through final token states, which is advantageous for long-range aggregation required by survival outcomes \citep{dai2019style,zhou2021informer}. In our setting, forecasting aligns with near-term prediction, whereas survival analysis emphasizes longer-range temporal dependencies; this mismatch naturally yields the observed performance split. While not a one-to-one architectural proof, the pattern is consistent with prior evidence on horizon-specific strengths \citep{wang2025large}.

\subsection{Interpreting very low $F_{1}$ scores and external validation drops}
 \textbf{Why some $F_{1}$ scores are near zero.} The very low F1 scores observed for certain models are not random failures but arise from two interacting factors. First, F1 is highly sensitive to threshold choice when the positive class is rare; small shifts in threshold calibration can yield 0.000 even when the model produces meaningful probabilities. Second, **architectural limitations** amplify this effect: **autoregressive decoders** tend to reproduce narrative rather than temporal order, and **instruction-tuned models** often default to generic or high-probability completions rather than selecting the correct event token—a behavior documented in prior analyses of instruction-following LLMs (\cite{liu2024lost, wei2021finetuned}).

These failures do not undermine model comparisons; rather, they highlight systematic architectural differences. **Encoder-based models**, trained with a **bidirectional masking objective** (\cite{devlin2019bert}), are better aligned with time-structured prediction tasks, as their contextual representation naturally integrates information from both past and future tokens. In contrast, **autoregressive decoders**, optimized for left-to-right generation (\cite{brown2020language}), are disadvantaged for temporal reasoning without explicit sequence restructuring or supervision. The contrast underscores that forecasting from clinical narratives is fundamentally different from standard text generation.
 
 \textbf{Why external validation is lower.} Performance drops on external subsets (e.g., \textit{sepsis-10} vs.\ \textit{sepsis-100}) reflect shifts in writing style and case mix, not random overfitting. Crucially, the \emph{relative ranking} of models remains stable across datasets, so comparative conclusions hold even when absolute scores vary.

% \subsection{Practical implications}
%  (i) Treat thresholding as part of the model: calibrate on a validation set (e.g., maximize $F_{1}$ or expected utility) and report threshold sensitivity. (ii) For decoders, impose explicit temporal supervision (sequence restructuring, time-aware heads) to mitigate left-to-right biases. (iii) Maintain cross-dataset evaluation to verify that relative rankings persist under domain/style shifts.


%%%%%%%%%%%%%%%%%%%%%%%%%%%%%%%%%%%%%%%%%%%%%%%%
%%%%%%%%%Appendix I: 
%%%%%%%%%%%%%%%%%%%%%%%%%%%%%%%%%%%%%%%%%%%%%%%
\section{Open-Source Models for Avoiding Potential Data Leakage}
\label{apd:open_source}

We evaluated potential data leakage arising from large language models (LLMs) pretrained on web-scale corpora that may include PubMed Open Access (PMOA) case reports. Our safeguards target two avenues of leakage: (i) \emph{content overlap} with PMOA text, and (ii) \emph{task leakage} from using representations too close to raw prose. Below we summarize the controls embedded in our methods and the additional empirical evidence gathered for this concern.

\subsection{Task Design Minimizes Leakage}

\textbf{Temporal annotations—not raw text—drive all evaluations.}
We transform narratives into chronological lists of \emph{(event, time)} tuples and operate on these derived annotations (future-masked where applicable). Models are never scored on lexical similarity to the raw prose.

\textbf{Temporal reordering changes the learning problem.}
Events are reordered by extracted timestamps, breaking the original textual sequence seen during pretraining. This creates a distinct modeling challenge: correct temporal ordering and causal dependencies must be inferred from structured tuples rather than reproduced from prose. As shown in the main text, training/evaluating on time-ordered versus text-ordered sequences yields measurably different performance.

\textbf{Temporal metrics reduce any advantage from memorization.}
Evaluation emphasizes temporal reasoning—forecasting F1 at fixed horizons, pairwise concordance (c-index), and AULTC—rather than surface-form similarity. Memorizing text order can even \emph{hurt} time-alignment metrics because clinical narratives often interleave retrospective commentary where textual order $\neq$ chronological order.

\subsection{Open-Source, PubMed-Free Baselines}

\textbf{Decoder baseline.} 
\emph{RedPajama-INCITE-7B-Instruct} is pretrained on English CommonCrawl (67\%), C4 (15\%), GitHub (4.5\%), Wikipedia (4.5\%), Books (4.5\%), arXiv (2.5\%), and StackExchange (2\%), explicitly \emph{excluding} PubMed. In our experiments it performs comparably to decoders that may have seen PMOA during pretraining.

\textbf{Encoder baselines.} 
We include \emph{BERT} (BooksCorpus + Wikipedia), \emph{RoBERTa} (CommonCrawl News, OpenWebText, Stories, Wikipedia), and \emph{DeBERTa} (\textit{RoBERTa}-style corpora). None of these encoders include PubMed in their pretraining data.


\subsection{Post-Cutoff External Validation (Additional Evidence)}
To probe encoder leakage further, we evaluated \textit{BERT}, \textit{RoBERTa}, and \textit{DeBERTa} on 115 case reports published \emph{after} their pretraining cutoffs (post-2020). Results mirror submission trends without degradation, which is inconsistent with a memorization hypothesis.

\begin{table}[!htbp]
\centering
\footnotesize
\setlength{\tabcolsep}{6pt}
\renewcommand{\arraystretch}{1.1}
\begin{tabular}{lccc}
\hline
\textbf{Model} & \textbf{0h} & \textbf{24h} & \textbf{168h} \\
\hline
BERT     & 0.63 & 0.62 & 0.51 \\
RoBERTa  & 0.54 & 0.61 & 0.65 \\
DeBERTa  & 0.58 & 0.56 & 0.53 \\
\hline
\end{tabular}
\caption{Post-cutoff evaluation on 115 case reports (forecasting F1 at 0h/24h/168h).}
\label{tab:postcutoff_f1}
\end{table}

We also report competitive performance from \emph{OLMO-32B-Instruct} (fully documented training data) and parity from the PubMed-excluded \emph{RedPajama-INCITE-7B-Instruct} relative to decoders potentially exposed to PMOA, further indicating that prior PubMed exposure does not meaningfully advantage our temporally structured tasks.

The \textbf{key takeaways} are 
\begin{itemize}
    \item Using derived, time-ordered \emph{(event, time)} tuples and temporal metrics directly targets temporal reasoning and reduces susceptibility to content memorization.
    \item PubMed-free baselines and post-cutoff tests provide external checks
    \item Empirically, models without PubMed exposure achieve comparable performance, arguing against causal leakage into our forecasting, concordance, and AULTC evaluations.
\end{itemize}


%%%%%%%%%%%%%%%%%%%%%%%%%%%%%%%%%%%%%%%%%%%%%%%%
%%%%%%%%%Appendix J:  
%%%%%%%%%%%%%%%%%%%%%%%%%%%%%%%%%%%%%%%%%%%%%%%
% \clearpage
\section{Generalizability Beyond Case Reports: Comparison with other Clinical Sources}
\label{apd:beyond_pmoa}

\subsection{Scope and Rationale for using Case Reports}
We acknowledge that published case reports differ from routine clinical documentation: they emphasize atypical presentations and are written for publication rather than point-of-care use, so they cannot substitute for validation on EHR notes. Our choice is methodological and pragmatic. Case reports are (i) publicly available, (ii) rich in temporally explicit clinical reasoning, and (iii) sufficiently structured to support extraction of textual time series. This makes them a practical, reproducible testbed for assessing whether forecasting models operate on temporally structured narratives without IRB barriers or restricted data access. We therefore use case reports to establish methodological feasibility and transparency, while recognizing that generalization to clinical operations must be demonstrated on EHR notes.

\subsection{Position Relative to Structured Clinical Scores}
Direct comparison to structured risk scores (e.g., SOFA, SAPS) is infeasible because required variables are absent from case reports. Our contribution is complementary to prior work that integrates text with structured data \citep{kline2022multimodal,wang2025large}: we show that free-text narratives, once transformed into timelines, carry prognostic signal usable for forecasting. Whereas \citet{kline2022multimodal} demonstrated that text can \emph{augment} structured models, our experiments isolate the incremental value of text when converted into time series, independent of structured covariates. In contexts where structured variables are available, these representations can be fused with established scores to assess additive value; here we focus on the text-only signal because the necessary structured inputs are not present in case reports.

%%%%%%%%%%%%%%%%%%%%%%%%%%%%%%%%%%%%%%%%%%%%%%%%
%%%%%%%%%Appendix K:  
%%%%%%%%%%%%%%%%%%%%%%%%%%%%%%%%%%%%%%%%%%%%%%%
\newpage
\section{Broader Societal Impact}
\label{apd:broader_impact}

This work advances equitable, scalable, and transparent medical AI by introducing a framework for forecasting clinical risk in sepsis from narrative documentation. Sepsis remains a leading cause of preventable mortality worldwide, with outcomes highly dependent on timely recognition and intervention. In many under‑resourced healthcare settings, structured EHR data and real‑time specialist input are unavailable, leaving unstructured narratives as the primary record of patient care. Our approach transforms such narratives into temporally structured representations, enabling systematic evaluation of forecasting methods for short‑ and long‑horizon event prediction, event ordering, and survival analysis. By releasing the dataset, extraction pipeline, and evaluation framework, we lower barriers for developing and benchmarking temporal forecasting models in sepsis, and promote reproducibility in LLM‑based clinical information extraction.

However, societal considerations are critical. LLM‑assisted extraction can introduce inaccuracies or hallucinations, and unvalidated forecasts risk amplifying bias or error. Although timelines are derived from non‑identifiable public case reports, their interpretation could still reinforce historical biases if not carefully contextualized. Forecasting results from case reports may also differ from those using routine clinical notes, limiting generalizability. Moreover, applying these forecasting models in high‑stakes clinical decision‑making without expert oversight poses safety risks. Responsible use requires transparent communication of dataset and model limitations, explicit uncertainty reporting, and safeguards to prevent inappropriate deployment in patient‑facing settings.

%\end{document}